\scopeparagraph{Software Engineering as our Target Discipline}

We target SE research because reporting guidelines from other disciplines do not address its specific needs (see \emph{Related Reporting Guidelines} below). In particular, SE research employs a wide variety of empirical methods~\cite{ralph2021empiricalstandardssoftwareengineering, DBLP:books/sp/WohlinRHORW24}. We organize the LLM-involving subset of these methods into a taxonomy of \studytypes\ifpaper~(Section~\ref{sec:study-types})\fi, and each of our \guidelines specifies its applicability per study type.

\scopeparagraph{Related Reporting Guidelines}

Reporting guidelines have a long tradition in healthcare research, with CONSORT~\cite{Schulz2010} (2010) as the canonical standard for randomized clinical trials. Our reporting checklist (\ifpaper Appendix~\ref{sec:checklist} \else \href{/checklist/}{checklist} \fi) follows its structural template. LLM-specific reporting guidelines have appeared more recently across communities: \citeauthor{Gallifant2025}'s TRIPOD-LLM~\cite{Gallifant2025} for healthcare, \citeauthor{DBLP:conf/chi/NavarroSA26}'s HCI guidelines~\cite{DBLP:conf/chi/NavarroSA26}, \citeauthor{Kapoor2024REFORMS}'s REFORMS~\cite{Kapoor2024REFORMS} for machine-learning-based science, and within SE, \citeauthor{DBLP:conf/icse/SallouDP24}'s vision paper on validity threats~\cite{DBLP:conf/icse/SallouDP24}, our earlier workshop position paper~\cite{DBLP:conf/wsese/0001BFB25}, and \citeauthor{DBLP:journals/corr/abs-2601-01954}'s prompt reporting guideline for automated SE~\cite{DBLP:journals/corr/abs-2601-01954}, derived from a literature review of ICSE, FSE, and ASE papers and a survey of 105 program committee members.

\citeauthor{DBLP:journals/corr/abs-2601-01954} frame their recommendations as complementing ours, with the essential items their 105 SE PC respondents endorsed lining up with what \modelversion, \design, and \limitationsmitigations already require. \citeauthor{DBLP:conf/chi/NavarroSA26}'s HCI guidelines, by contrast, balance \enq{the practical realities of authors' time, cost, and page limitations} with scientific concerns and HCI norms. Our \guidelines instead aim for comprehensive coverage and use \must and \should within each to express the balance. On prompt reporting, \citeauthor{DBLP:conf/chi/NavarroSA26} argue for selective reporting based on each prompt's centrality to author claims, whereas \design recommends fuller prompt disclosure with named exceptions for privacy, anonymity, or confidentiality concerns. Peripheral LLM uses such as proof-reading, spell-checking, and translation are out of scope here entirely (see \emph{Focus on Direct Development or Research Support} below). On technical evaluation, \citeauthor{DBLP:conf/chi/NavarroSA26} recommend a \enq{modest technical evaluation of the LLM component on a dataset of representative inputs} tailored to HCI research, whereas \benchmarksmetrics asks for established benchmarks, traditional baselines, and inferential statistics in the study types where that applies.

Our guidelines name \paper or \supplementarymaterial as the reporting location for each recommendation (see \emph{Reporting Locations} below). The \supplementarymaterial is published according to the \emph{ACM SIGSOFT Open Science Policies}~\cite{daniel_graziotin_2024_10796477} for replication and downstream reuse. In summary, our guidelines apply to SE empirical research broadly, organize recommendations around the taxonomy of seven study types introduced in \studytypes with explicit per-study-type applicability, and pair each recommendation with a target reporting location in \paper or \supplementarymaterial.

\scopeparagraph{Focus on Text-Based Use Cases}

While multi-modal foundation models that use or generate images, audio, or video may also support SE research and practice, we focus on textual use cases of LLMs (e.g., in natural language or programming languages). Many of our guidelines, particularly those concerning model reporting, prompt documentation, and reproducibility, are likely applicable to multi-modal settings as well, but we leave their explicit validation to future work.

\scopeparagraph{Focus on Direct Development or Research Support}

While researchers may use LLMs for many peripheral tasks (e.g., proof-reading, spell-checking, translation), our guidelines focus on their direct role in empirical research and engineering practice.
For engineers, we focus on the use of LLMs to automate SE tasks, that is, artificial intelligence (AI) for software engineering (AI4SE) (see~\llmsforengineers).
This includes agentic systems that autonomously plan and execute multi-step tasks using LLMs (see~\design).
For researchers, we focus on the use of LLMs to automate empirical research tasks such as data collection, processing, or analysis (see~\llmsforresearcher).

\scopeparagraph{Researchers as our Target Audience}

Our guidelines are intended to help SE researchers design, plan, conduct,
and report empirical studies involving LLMs, and to support scholarly peer
review of such studies. Each guideline includes an \emph{Advice for
Reviewers} subsection with targeted assessment suggestions.
Our guidelines focus on what to report and how; they complement but do not replace methodological guidance for designing specific types of empirical studies.

\scopeparagraph{Reporting Locations}

We use \paper to refer to the manuscript PDF, including any appendices bound with it. We use \supplementarymaterial to refer to any artifact, replication package, dataset, or other resource external to the manuscript PDF (e.g., hosted on Zenodo or Figshare). What belongs in the main body versus an appendix is a presentation choice for authors and venues; what belongs inside the manuscript versus outside it is the reporting-location decision the guidelines make. When a recommendation does not specify a location, either \paper or \supplementarymaterial is acceptable.

\scopeparagraph{Navigating These Guidelines}

These guidelines are structured to support different reading strategies depending on the reader's goal. \emph{Researchers planning a new study} may start with the taxonomy of study types (see \studytypes) to identify which types apply to their planned work, then consult \ifpaper Table~\ref{tab:guideline-matrix} \else the \href{/guidelines/\#guidelines-by-study-type}{applicability matrix} \fi to determine which guidelines are requirements (\must) and which are recommendations (\should) for those study types. Each guideline section opens with a brief summary in a shaded box, allowing readers to quickly assess relevance before reading the full text. \emph{Researchers writing up results} may prefer to start with the \ifpaper checklist in Appendix~\ref{sec:checklist}, \else \href{/checklist/}{checklist}, \fi which organizes actionable items by typical paper sections (Introduction, Research Design and Methods, Results, etc.). \ifpaper Table~\ref{tab:guidelines} provides a quick reference to all eight guidelines, and Table~\ref{tab:rationale-recommendations} maps each guideline's rationale to its recommendations. \fi \emph{Reviewers} can use the \emph{Advice for Reviewers} subsection at the end of each guideline for targeted guidance on assessing manuscripts. \ifpaper Table~\ref{tab:guideline-matrix} helps reviewers identify which guidelines apply to the study type under review. \fi